\section{Statements and scope}\label{the-complete-principal-contract}

Hou, Feng and Leemans {[}HFL{]} study regular polyhedra whose automorphism
groups have order \(2^n p\), where \(p\) is an odd prime. This application
recovers their principal type restriction, both families of cubic
constructions, their general odd-factor construction, and the exact
existence range for type \(\{6,2^s\}\). We use their arXiv:2001.02945v2
edition, dated 14 July 2024, comprising seventeen pages. The correction
in Section 6 concerns an auxiliary presentation in that edition; no
claim about the later journal display is made.

We work with marked rank-three string C-groups. The distinguished
generators \(a,b,c\) are involutions, \(ac=ca\), and
\(\langle a,b\rangle\cap\langle b,c\rangle=\langle b\rangle\).
This last equality is the complete rank-three intersection test. We
retain the standard correspondence with regular abstract polyhedra,
under which the group acts regularly on flags and its order is the
number of flags. Nondegeneracy means both adjacent product orders are
at least three. Our conventions are \(x^g=g^{-1}xg\) and \([x,y]=x^{-1}y^{-1}xy\).

\textbf{Theorem 1 (type restriction; {[}HFL, Theorems 1.2 and 3.1{]}).} If a
group of order \(2^n p\) is generated by three elements of order at most two with commuting
endpoints, then \(p\) divides at least one adjacent product order.
Consequently, up to duality, a nondegenerate regular polyhedron of this
order has type \(\{2^s p,2^t\}\) or \(\{2^s p,2^t p\}\).

\textbf{Theorem 2 (odd-factor extension; {[}HFL, Theorem 4.1{]}).} For positive
odd integers \(\ell_1,\ell_2\), integers \(s,t\geq2\) and
\(n\geq s+t+1\), the explicit presentation in Section 3 is a solvable
string C-group of type \(\{2^s\ell_1,2^t\ell_2\}\) and order
\(2^n\ell_1\ell_2\). In particular it gives the source's Corollary 4.2
by taking \(\ell_1=p\), \(\ell_2=1\).

\textbf{Theorem 3 (the exact six-sided range; {[}HFL, Theorem 1.3{]}).} For
\(n\geq7\), a nondegenerate regular polyhedron of order \(3\cdot2^n\)
and type \(\{6,2^s\}\) exists if and only if

\[2\leq s\leq n-2,\qquad s\ne n-3.\tag{1}\]

The constructions identify the source's actual presentations, with the
single sign correction specified in Section 6.

\textbf{Theorem 4 (the cubic families; {[}HFL, Theorems 4.3 and 5.1{]}).} The
three literal presentations of type \(\{6,4\}\) give solvable groups of
orders \(96m^3,192m^3,384m^3\), for every integer \(m\geq1\). The source's
three type-\(\{6,6\}\) presentations give solvable groups of orders
\(192m^3,384m^3,768m^3\), for every \(m\geq1\). In particular, type
\(\{6,4\}\) and type \(\{6,6\}\) both occur at order \(3\cdot2^n\)
for every \(n\geq5\).

The shared ingredients are the supported-kernel and square-group
results included in [U], and its complete affine-lattice criterion.
The existing exact \(\{6,6\}\) calculation {[}C66{]} is reused. Burnside's
\(p^a q^b\) solvability theorem and Lucchini's cyclic-core bound are
retained classical inputs, as in the source. One bounded exclusion
retains the complete GAP SmallGroups catalogue of order 384. These
inputs, rather than experimental extrapolation, are the precise
boundaries of the recovery.
